\documentclass[11pt]{article}
\usepackage[margin=1in]{geometry}
\usepackage{amsmath,amssymb,amsthm}
\usepackage{hyperref}

\newcommand{\query}{\mathcal{Q}}
\newcommand{\block}{F}
\newcommand{\interface}{I}
\newcommand{\joinstate}{S}
\newcommand{\muLp}[3]{\mu_{#1}(#2 \mid #3)}
\newcommand{\attr}[1]{\mathrm{attr}(#1)}

\newtheorem{theorem}{Theorem}
\newtheorem{lemma}{Lemma}

\title{Statistics Before Programs: Block $\ell_p$-Norms for Cyclic Join Bounds\\[4pt]
\large Extended version: full proofs}
\author{Suan Lee\\ Semyung University\\ \texttt{suanlee@semyung.ac.kr}}
\date{}

\begin{document}
\maketitle

\noindent This note accompanies the ICDE 2027 submission and gives the
proofs that the paper states in condensed form: the $\ell_p$ Composition
Inequality (the validity of every LpH-CEG bound) and the Cyclic
Improvement theorem. Notation follows the paper.

\section{The $\ell_p$ Composition Inequality}

\paragraph{Setup.} A hyperpath $P=(\joinstate_0{=}\emptyset,\block_1,
\joinstate_1,\dots,\block_m,\joinstate_m)$ adds a block of relations
$\block_j$ at step $j$, $\joinstate_j=\joinstate_{j-1}\cup\block_j$, and
$\joinstate_m$ is the set of all relations. Write $A_j=\attr{\block_j}$,
$V_j=\attr{\joinstate_j}$ ($V_0=\emptyset$, $V_m=V$), the interface
$\interface_j=A_j\cap V_{j-1}$ and $E_j=A_j\setminus\interface_j$. Then
\begin{itemize}
\item[(F1)] $E_j=V_j\setminus V_{j-1}$, so $\{E_j\}_{j}$ partitions $V$;
\item[(F2)] $\interface_j\subseteq V_{j-1}$.
\end{itemize}
Weights $w_{j,p}\ge0$ ($p\in\mathcal{P}\subseteq[1,\infty]$) satisfy the
attribute-wise cover constraint of the paper: for every attribute $a\in V$,
bound at step $j(a)$ (i.e.\ $a\in E_{j(a)}$),
\begin{equation}\tag{C}
w_{j(a)}+\sum_{j':\,a\in\interface_{j'}}s_{j'}\;\ge\;1,
\qquad w_j:=\sum_p w_{j,p},\quad s_j:=\sum_p w_{j,p}/p\ \ (1/\infty=0).
\end{equation}
A \emph{polymatroid} is $h:2^V\to\mathbb{R}$ with $h(\emptyset)=0$,
monotone and submodular; $h(X\mid Y):=h(X\cup Y)-h(Y)$.

\begin{lemma}[Conditioning]\label{lem:cond}
For a polymatroid $h$ and $Y\subseteq Z$: $h(X\mid Y)\ge h(X\mid Z)$.
\end{lemma}
\begin{proof}
Submodularity on $X\cup Y$ and $Z$ gives $h(X\cup Y)+h(Z)\ge
h(X\cup Y\cup Z)+h((X\cup Y)\cap Z)$, and $h((X\cup Y)\cap Z)\ge h(Y)$ by
monotonicity since $Y\subseteq(X\cup Y)\cap Z$. Rearrange.
\end{proof}

\begin{lemma}[Fractional subadditivity]\label{lem:frac}
Let $f$ be monotone submodular on subsets of a finite set with
$f(\emptyset)=0$, and let $\lambda_i\ge0$ satisfy
$\sum_{i:\,a\in T_i}\lambda_i\ge1$ for every $a\in S$. Then
$\sum_i\lambda_if(T_i)\ge f(S)$.
\end{lemma}
\begin{proof}
Order $S=(a_1,\dots,a_k)$ and let $S_{<t}=\{a_1,\dots,a_{t-1}\}$. By
monotonicity, the chain rule and Lemma~\ref{lem:cond},
$f(T_i)\ge f(T_i\cap S)=\sum_{a_t\in T_i}f(a_t\mid T_i\cap S_{<t})
\ge\sum_{a_t\in T_i}f(a_t\mid S_{<t})$. Weighting by $\lambda_i$ and
exchanging sums, $\sum_i\lambda_if(T_i)\ge\sum_t f(a_t\mid S_{<t})
\sum_{i:\,a_t\in T_i}\lambda_i\ge\sum_t f(a_t\mid S_{<t})=f(S)$, using
$f(a_t\mid S_{<t})\ge0$.
\end{proof}

\begin{lemma}[The obligation]\label{lem:obl}
For every polymatroid $h$ and every $w$ satisfying (C),
\[
\sum_j\Bigl[s_j\,h(\interface_j)+w_j\,h(E_j\mid\interface_j)\Bigr]\;\ge\;h(V).
\]
\end{lemma}
\begin{proof}
Let $g_j=h(E_j\mid V_{j-1})$, $\delta_j=\max(0,1-w_j)\in[0,1]$ and
$f_j(S)=h(S\mid V_{j-1})$ for $S\subseteq E_j$; $f_j$ is monotone,
submodular and $f_j(\emptyset)=0$.

(i) For each $j'$, the chain rule along the partition (F1) and
Lemma~\ref{lem:cond} give
$h(\interface_{j'})=\sum_j h(\interface_{j'}\cap E_j\mid
\interface_{j'}\cap V_{j-1})\ge\sum_j f_j(\interface_{j'}\cap E_j)$.
Hence $\sum_{j'}s_{j'}h(\interface_{j'})\ge\sum_j\sum_{j'}s_{j'}
f_j(\interface_{j'}\cap E_j)$.

(ii) Fix $j$. If $\delta_j>0$ then by (C) every $a\in E_j$ has
$\sum_{j':a\in\interface_{j'}}s_{j'}\ge1-w_j=\delta_j$, so the sets
$\interface_{j'}\cap E_j$ with weights $s_{j'}/\delta_j$ fractionally
cover $E_j$, and Lemma~\ref{lem:frac} gives
$\sum_{j'}s_{j'}f_j(\interface_{j'}\cap E_j)\ge\delta_j\,g_j$. If
$\delta_j=0$ the same inequality holds because the left side is
non-negative.

(iii) By (F2) and Lemma~\ref{lem:cond}, $h(E_j\mid\interface_j)\ge g_j\ge0$,
and $w_j\ge1-\delta_j$, so $w_jh(E_j\mid\interface_j)\ge(1-\delta_j)g_j$.

Adding (i)--(iii): the left side is at least
$\sum_j[\delta_j+(1-\delta_j)]g_j=\sum_jh(E_j\mid V_{j-1})=h(V)$ by the
chain rule and (F1).
\end{proof}

\begin{theorem}[$\ell_p$ Composition Inequality]
Let $D$ be a database, $\query(D)\ne\emptyset$, and let
$\muLp{p}{\block_j}{\interface_j}$ be the $\ell_p$-norm of the degree
sequence of $E_j$-values per $\interface_j$-value in a relation
$R_j\supseteq\pi_{A_j}\query(D)$---for instance the block join
$\bowtie_{R\in\block_j}R$ computed exactly on $D$. If $w$ satisfies (C), then
\[
|\query(D)|\;\le\;\prod_{j=1}^m\prod_{p\in\mathcal P}
\muLp{p}{\block_j}{\interface_j}^{\,w_{j,p}}.
\]
\end{theorem}
\begin{proof}
Let $h$ be the entropy of the uniform distribution on $\query(D)$; it is a
polymatroid and $h(V)=\log_2|\query(D)|$. The marginal on $A_j$ is
supported on $\pi_{A_j}\query(D)\subseteq R_j$, so the $\ell_p$ degree
inequality of Abo Khamis et al.\ (PODS 2024) gives, for $p\in[1,\infty]$,
$\frac1p h(\interface_j)+h(E_j\mid\interface_j)\le
\log_2\muLp{p}{\block_j}{\interface_j}$. Multiply by $w_{j,p}\ge0$, sum, and
apply Lemma~\ref{lem:obl}.
\end{proof}

\paragraph{Scope.} The proof uses only monotonicity and submodularity, so
it holds over the whole polymatroid cone and for every hyperpath, block
size and $p\ge1$. It requires (C) exactly: an LP solver may return weights
violating (C) by about $10^{-7}$, which is removed by scaling $w$ by
$1/\min_a(\text{coverage of }a)$. It requires exact statistics of the
database being bounded (or deterministic upper bounds on them, since the
bound is monotone in each $\mu$); a sampled or sketched statistic, such as
the Hutchinson estimate of the common-neighbour $\ell_2$ in the paper, is
outside the guarantee. When some block join is empty the query is empty
and the bound is $0$.

\paragraph{Why the elementary argument does not suffice.} If $w_j\ge1$ for
every step, Lemma~\ref{lem:obl} follows from (iii) alone. The LP's optima
on cyclic queries have $w_j<1$ (e.g.\ $1/3$ on the closing step of a
cycle); the deficit $\delta_j$ is paid for by later interfaces, which is
exactly step (ii). The cyclic improvement and the need for
Lemma~\ref{lem:frac} have the same source.

\paragraph{Independent checks.} An LP over the polymatroid cone decides
Lemma~\ref{lem:obl} at fixed weights; enumerating every vertex of the
feasible weight region decides it for all weights. Over 6{,}197 hyperpath
structures on eleven shapes with $|V|\le6$ and block sizes $1$--$3$ the
minimum slack is exactly $0$; an adversarial search over random abstract
hyperpath schemes with $|V|\le5$ found no violation; and over $207{,}000$
databases the bound never fell below the exact count
(\texttt{scripts/shannon\_check\_wide.py},
\texttt{scripts/validity\_search.py}, \texttt{docs/step5\_lemma.md}).

\section{The Cyclic Improvement theorem}

Order the cycle $C_k$ as $R_1(a_1,a_2),\dots,R_k(a_k,a_1)$ and extend in
this order with single-relation blocks. Let $N_j=|R_j|$, $M_j$ the maximum
degree and $\mu_p(R_j)$ the $\ell_p$-norm of the degree sequence of $R_j$
from $a_j$ to $a_{j+1}$.

\paragraph{Structure.} Step 1 has interface $\emptyset$ and extension
$\{a_1,a_2\}$: its degree sequence is the single value $N_1$. A middle step
$2\le j\le k-1$ has interface $\{a_j\}$, extension $\{a_{j+1}\}$ and
$\mu_\infty=M_j$. The closing step $k$ has interface $\{a_k,a_1\}$ and
empty extension; its conditional degrees are all $1$, so $\mu_\infty=1$
and $\mu_1=N_k$. With $p=\infty$ and unit weights this gives
$B_{\max}=N_1\prod_{j=2}^{k-1}M_j$.

\paragraph{A feasible point for every common $p\ge1$.} Put weight
$\frac{p}{p+1}$ with exponent $p$ on every middle step, $\frac{p}{p+1}$ on
step 1 and $\frac1{p+1}$ on the closing step with exponent $1$. The cover
constraints read: $a_1$: $\frac p{p+1}+\frac1{p+1}=1$; $a_2$:
$\frac p{p+1}+\frac{p/(p+1)}{p}=1$; $a_i$ for $3\le i\le k-1$:
$\frac p{p+1}+\frac1{p+1}=1$; $a_k$: $\frac p{p+1}+\frac1{p+1}=1$. All are
tight, so the LP optimum satisfies
\[
B_{\mathrm{LP}}\le N_1^{\frac p{p+1}}N_k^{\frac1{p+1}}
\prod_{j=2}^{k-1}\mu_p(R_j)^{\frac p{p+1}}
\quad\text{and}\quad
\frac{B_{\max}}{B_{\mathrm{LP}}}\ge\Bigl(\frac{N_1}{N_k}\Bigr)^{\frac1{p+1}}
\prod_{j=2}^{k-1}\frac{M_j}{\mu_p(R_j)^{p/(p+1)}}.
\]
\paragraph{The exponent must be common (or non-increasing).} The
constraint of the attribute between middle edges $j$ and $j{+}1$ is
$\frac{p_j}{p_j+1}+\frac1{p_{j+1}+1}\ge1$, i.e.\ $p_{j+1}\le p_j$. Choosing
each edge's best exponent independently is therefore not feasible in
general (e.g.\ $p_j=1$, $p_{j+1}=\infty$ covers that attribute only to
$\frac12$); the LP may still mix exponents within a step.

\paragraph{Growth.} At $p=2$ a middle factor $M_j/\mu_2(R_j)^{2/3}$ exceeds
$1$ exactly when $\mu_2(R_j)<M_j^{3/2}$, i.e.\ for skewed relations, so
under fixed skew the ratio grows geometrically in the number $k-2$ of
middle edges. As $p\to\infty$ every factor tends to $1$, and LpH-CEG
returns $\min(B_{\max},B_{\mathrm{LP}})$, so the realized ratio is always
at least $1$. The boundary factor $(N_1/N_k)^{1/(p+1)}$ does not depend on
$k$; it is $1$ for a uniform cycle.

\paragraph{Acyclic queries.} On a star the cover constraints force unit
weight on every step and the LP returns the max-degree bound. On a path the
LP can use an interface at $p=1$ in place of an extension (on a 4-path it
returns weights $(1,0,1)$), so acyclicity alone does not force unit
weights; what redundant coverage on cycles adds is the slack exploited
above.

\section{Complexity}
Collecting $\muLp{p}{\block}{\interface}$ for a single relation is one
scan, $O(N)$; for a two-relation block it materializes the pair join,
except for the interfaces computable from single-relation degree tables
(all but the common-neighbour interface of a binary block; Section VI-C of
the paper). The cover LP has $m|\mathcal P|$ variables and $|V|$
constraints; it solves in milliseconds for the queries in the paper. The
flow LP over a $bs\le2$ catalog is polynomial in the number of constraints
and solves in $0.03$--$4.4$\,s at $k=3$--$16$ over $4k$ orders.

\end{document}
